%! Concluding a Test for a Population Proportion — Unit 6, Lesson 6.6 — Homework
\documentclass[10pt]{article}
\usepackage{apstats-article}
\usepackage{apstats-boxes}
\newcommand{\TallMath}[1]{$\displaystyle #1\rule[-1.4em]{0pt}{3.2em}$}

\begin{document}

\pageheader{Unit 6, Lesson 6.6}{Homework}
\namedateperiod

\vspace{0.15in}

\textit{Perform a complete four-step hypothesis test (State, Plan, Do, Conclude) for each problem.}

\begin{enumerate}

  \item A marketing firm claims 40\% of adults in a city have visited a particular
        website in the past month. A competitor believes the proportion is \emph{lower}.
        From a random sample of 350 adults, 126 have visited the site. Use $\alpha=0.05$.

        \textbf{State:} \writelines{1}
        \textbf{Plan:} \writelines{2}
        \textbf{Do:} $\hat p = $ \blank{2cm} \quad $z \approx$ \blank{2cm} \quad $p$-value $\approx$ \blank{2cm}
        \textbf{Conclude:} \writelines{2}

  \item A national study reports that 72\% of high school students participate in at
        least one extracurricular activity. A district administrator believes the rate
        in her district has \emph{changed}. From 250 randomly selected students, 193
        participate. Use $\alpha=0.05$.

        \textbf{State:} \writelines{1}
        \textbf{Plan:} \writelines{2}
        \textbf{Do:} $\hat p = $ \blank{2cm} \quad $z \approx$ \blank{2cm} \quad $p$-value $\approx$ \blank{2cm}
        \textbf{Conclude:} \writelines{2}

  \item A bottled water company claims fewer than 5\% of their bottles have a
        contamination level above the safe threshold. An environmental group tests 400
        randomly selected bottles; 26 exceed the threshold. Use $\alpha=0.01$.

        \textbf{State:} \writelines{1}
        \textbf{Plan:} \writelines{2}
        \textbf{Do:} $\hat p = $ \blank{2cm} \quad $z \approx$ \blank{2cm} \quad $p$-value $\approx$ \blank{2cm}
        \textbf{Conclude:} \writelines{2}

\end{enumerate}

\begin{reflectionbox}
\textbf{Preview:} Lesson 6.7 introduces Type I errors, Type II errors, and power. Reflect:
when you fail to reject $H_0$, does that mean $H_0$ is definitely true? What might be
happening if you keep getting $p > \alpha$ even when $H_a$ is actually true?
\end{reflectionbox}

\end{document}
